\section{Conclusion}\label{sec:conclusion}

We set out to build a text-to-motion model that treats the skeleton as an input rather than as a
part of the architecture: a per-joint representation any rigged animation converts to, normalised
cell by cell; a diffusion transformer that reads a rest frame, joint descriptions and the kinematic
tree and denoises the raw representation; and a grouped objective calibrated from measured
energies. One \num{303}M-parameter model trained on \num{311} animal rigs generates text-conditioned
motion for all of them (R@1 \num{0.922} on held-out clips, against \num{0.964} for the real
clips), and a rig it has never seen is served by the same forward pass and adapted by a low-rank adapter trained on eight to twenty-six
clips that brings its rotation and position channels back into agreement.

\paragraph{Limitations.} The library is animal-only. Retrieval metrics come from an evaluator we
train, hence the real-motion reference next to every number, and rigs outside the library are
scored geometrically. The extension study is four case studies with two to six
validation clips each, and language following outside an adapter's clips is measured on one rig (App.~\ref{app:extend_more}).

\section*{Reproducibility statement}

Section~\ref{sec:method:rep} and Appendix~\ref{app:rep} specify the representation channel by
channel, the conversion from a kinematic tree, rest pose and motion, and the per-cell
normalisation; Appendix~\ref{app:calib} gives the calibration protocol, the target share profile
and the mechanism check that certifies an artifact; Appendix~\ref{app:lora} lists every adapter
hyper-parameter; Appendix~\ref{app:anytop} gives the AnyTop sampling command and the geometric
protocol, including the frame-rate handling. All evaluation runs are seeded, and the sharded
generation used for retrieval metrics records the checkpoint, sampler and source hashes so that
a table can be re-scored from the stored samples.

\subsection*{AI use statement}
We used generative AI tools throughout this work. The code for data conversion, training, evaluation
and rendering was drafted with an LLM-based coding assistant and independently reviewed by a second
LLM reviewer before it was run; every script was tested, and every number in this paper was checked
against the run logs and artifacts by the authors. LLM assistance was also used for experiment
orchestration, for auditing the data, for figure scripts, literature search
and the drafting and editing of the text. The natural-language captions of the corpus and the joint
descriptions were produced with an LLM from the animation metadata and joint names. The research
questions, the design decisions and the interpretation of the results are the authors', and we take
responsibility for the final content, including text, claims and artifacts produced with the aid of
generative AI.

\subsection*{Ethics statement}
The work involves no human subjects and no personal data. The training library is derived from the
animation assets of a commercial game and the evaluation rigs from the publicly distributed
Truebones zoo; the game assets are used for research and are not redistributed, and the paper
describes the conversion so that the representation can be rebuilt from licensed copies
(\todo{confirm the data-source wording with the authors}). Generated animal motion for games and film
carries no risk beyond the generic misuse of animation tools.
